\documentclass[11pt]{letter}

\usepackage[margin=0.78in]{geometry}
\usepackage[T1]{fontenc}
\usepackage[utf8]{inputenc}
\usepackage{lmodern}
\usepackage[hidelinks]{hyperref}
\usepackage{setspace}

\address{School of Aerospace\\
Harbin Institute of Technology\\
Shenzhen 518055, China\\
\href{mailto:liujian2024@hit.edu.cn}{liujian2024@hit.edu.cn};
\href{mailto:hitchenbo@hit.edu.cn}{hitchenbo@hit.edu.cn}}
\date{31 August 2026}

\begin{document}

\begin{letter}{Editors\\
\emph{Remote Sensing}}

\opening{Dear Editors of \emph{Remote Sensing},}

Please consider our manuscript entitled ``Statistical Evidence for a DQN-Family Advantage in a Resource-Coupled Satellite-Ground Scheduling Benchmark'' for publication as an Article in \emph{Remote Sensing}.

Satellite systems supporting Earth observation must allocate computation, energy, thermal capacity, bandwidth, and contact time across sequential workloads. Existing studies have advanced constellation scheduling, observation planning, service placement, and reinforcement-learning controllers. However, controlled evidence separating the benefit of long-horizon value learning from the choice of a particular deep Q-network (DQN) training objective remains limited in resource-coupled satellite-ground scheduling.

Our study establishes a fully specified synthetic benchmark and evaluates six DQN objectives under matched architectures, training budgets, model seeds, and held-out traces. Across all 30 objective--regime comparisons, every DQN objective achieved a lower mean episode cost than four-step model-predictive control (MPC-4). All paired 95\% bootstrap intervals excluded zero, and Holm-adjusted sign tests supported the same ordering. A correlation-preserving joint bootstrap and a simultaneous upper bound assess this family-level result without selecting only the best-performing variant. The manuscript also reports the smaller, scenario-dependent differences within the DQN family, the cost--runtime trade-off of deeper exact planning, sensitivity analyses, and robustness to training-time preview misspecification.

The manuscript fits the scope and readership of \emph{Remote Sensing} because it examines reproducible algorithms for satellite-ground processing and the management of persistent onboard and communication resources relevant to responsive remote-sensing operations. The conclusions are explicitly bounded to a synthetic benchmark with soft constraints, and the accompanying public reproducibility package provides a transparent basis for subsequent mission-calibrated, hard-constrained, and operational validation.

We confirm that neither the manuscript nor any parts of its content are currently under consideration for publication with or published in another journal.

All authors have approved the manuscript and agree with its submission to \emph{Remote Sensing}.

Thank you for your consideration.

\vspace{0.8em}
Sincerely,\\[1em]
Jian Liu and Bo Chen\\
Corresponding Authors, on behalf of all authors

\end{letter}
\end{document}
